\section{Consequences for deployment choices}
\label{app:decision}
The preceding comparisons measure classification and cost-prediction errors.
Here we retrospectively apply the same deployment-choice rule to each
campaign's recorded non-anchor candidate pool: require predicted P1+ in
\emph{every} window, then select the policy with the lowest mean predicted
window cost. Windows receive equal weight. Archived policy ID breaks exact
cost ties at the released data precision; all tied minimizers remain in the
artifact. Physical results do not enter this selection rule.

For evaluation, a selected policy is feasible only if every recorded physical
repeat passes P1+. Its observed cost first averages repeats within each window,
then averages windows equally. Let $c_*$ be the lowest such observed mean
among candidates feasible in all their runs. We report the cost gap
$100(c_{\rm selected}/c_*-1)$ only for feasible selections. This is a
comparison within the tested pool, not a global optimum; all2 capacity anchors
are excluded. Each campaign retains its own workload, runtime and price contract.

\begin{table*}[t]
\centering
\small
\caption{Retrospective policy choices from the original frozen predictions.
Costs are scenario USD; a dash excludes an infeasible selection from the
feasible-cost comparison. P1+ counts repeated physical executions, not independent
choices. $\dagger$ marks a prediction-cost tie between H and H+guard, resolved
by archived ID to H. E5's two tied policies are physically feasible; E6's H
fails whereas H+guard passes.}
\label{tab:decisionutility}
\par\vspace{4pt}
\begin{tabular*}{\linewidth}{@{\extracolsep{\fill}}llllrr@{}}
\toprule
Study & Replay & Selected policy & Physical P1+ & Observed cost & Feasible cost gap (\%)\\
\midrule
E2 & C & C30-2 & 6/6 & 2.0769 & 0.00 \\
E2 & C30 & C-2 & 0/6 & 2.8398 & -- \\
E3 & C & N1 & 4/4 & 1.9423 & 2.70 \\
E3 & C30 & N2 & 4/4 & 1.8913 & 0.00 \\
E4 & C & N2 & 4/4 & 1.8814 & 58.95 \\
E4 & C30 & H+guard & 4/4 & 1.1836 & 0.00 \\
E4 & C$_{\rm sg}$ & H+guard & 4/4 & 1.1836 & 0.00 \\
E5 & C & C-1 & 6/6 & 2.1655 & 94.13 \\
E5 & C30 & H+guard & 6/6 & 1.1155 & 0.00 \\
E5 & C$_{\rm sg}$ & H$^{\dagger}$ & 6/6 & 1.1189 & 0.31 \\
E6 & L14 & H+guard & 4/4 & 1.6019 & 0.00 \\
E6 & S14 & H$^{\dagger}$ & 0/4 & 1.5149 & -- \\\bottomrule

\end{tabular*}
\end{table*}

Table~\ref{tab:decisionutility} makes the practical consequences explicit.
E2's C30 has a unique cost-minimizing candidate that fails every physical
repeat. In E4 and E5, C excludes H+guard, the lowest-cost feasible policy
in those measured pools; its chosen policies cost 58.95\% and 94.13\% more,
respectively. C30 selects H+guard in both. These are opportunity costs under
the declared synthetic-API ledger, not commercial savings attributed to a
new optimizer. E5's C$_{\rm sg}$ ties H and H+guard in prediction; their
0.31\% observed cost difference is descriptive and does not establish a
reliable ranking. E6's S14 also ties these policies, so a failing choice is
possible under the stated rule rather than inevitable under every tie-break.

The candidate pools contain five, six, six, four and two policies in E2--E6,
respectively. E3 therefore compares its six physically evaluated checkpoints,
not all 90 screened candidates. This reanalysis adds no physical run or
prospective selection test and does not choose a universally preferred replay
model. The source bundle's \path{artifact/decision_utility/} reconstructs the
choices and table from reduced per-record metrics; the paired corpus provides
the underlying request and accounting records.


\subsection{Price robustness of deployment choices}
We reprice the same recorded events at GPU and synthetic-API price
multipliers in $\{0.5,1,2\}\times\{0.5,1,2\}$, keeping startup, shutdown,
and deadline-miss charges unchanged. At each point, we repeat the prediction-only
selection and recompute the lowest-cost physically feasible comparator within
the same pool. Table~\ref{tab:decisionprice} summarizes all 108 choices; ranges
are grid extrema, not confidence intervals or new executions.

\begin{table*}[t]
\centering
\small
\caption{Deployment-choice sensitivity over all nine price combinations.
Feasible counts refer to price combinations whose selected policy passes all
recorded physical repeats. $\dagger$ retains the H/H+guard prediction tie:
ID order selects H, but H+guard is also a minimizer.}
\label{tab:decisionprice}
\par\vspace{4pt}
\begin{tabular*}{\linewidth}{@{\extracolsep{\fill}}llllr@{}}
\toprule
Study & Replay & Selected policies & Feasible price cells & Cost gap range (\%)\\
\midrule
E2 & C & C30-2, C-1 & 9/9 & 0.00--11.04 \\
E2 & C30 & C-2 & 0/9 & -- \\
E3 & C & N2, N1 & 9/9 & 0.00--3.11 \\
E3 & C30 & N2, N1 & 9/9 & 0.00 \\
E4 & C & N2 & 9/9 & 16.56--80.17 \\
E4 & C30 & H+guard & 9/9 & 0.00 \\
E4 & C$_{\rm sg}$ & H+guard & 9/9 & 0.00 \\
E5 & C & C-1 & 9/9 & 41.93--119.65 \\
E5 & C30 & H+guard & 9/9 & 0.00 \\
E5 & C$_{\rm sg}$ & H$^{\dagger}$ & 9/9 & 0.27--0.33 \\
E6 & L14 & H+guard & 9/9 & 0.00 \\
E6 & S14 & H$^{\dagger}$ & 0/9 & -- \\\bottomrule

\end{tabular*}
\end{table*}

E2's C30 uniquely selects the failing C-2 throughout the grid. E4 and E5's C
selections remain feasible but costlier than the tested feasible minimum.
E2's C and both E3 predictors change their choices with prices, confirming that
the analysis reselects policies rather than simply repricing a fixed choice.
The E5 and E6 ties persist throughout: E5's small H/guard cost difference does
not establish a reliable ranking, and E6's failing choice depends on the
tie-break. The source bundle's \path{artifact/decision_price/} retains every
price point, source-bound cost component, and selection, and exactly recovers
the nominal-price results above. Fixed-event repricing does not model policies
retrained or redeployed under new prices.
